\DocumentMetadata{
  pdfversion=2.0,
  pdfstandard=ua-2,
  lang=en-US,
  tagging=on,
  tagging-setup={math/setup=mathml-SE,tabsorder=structure}
}
\documentclass[11pt,letterpaper]{article}
\usepackage[margin=0.68in,includefoot]{geometry}
\usepackage{newcomputermodern}
\usepackage{amsmath,amscd,mathtools}
\usepackage{tikz,adjustbox}
\providecommand{\square}{\mdlgwhtsquare}
\usepackage[hidelinks]{hyperref}
\hypersetup{
  pdftitle={Analysis PhD Exam, January 1997},
  pdfauthor={University of Florida Department of Mathematics},
  pdfsubject={Graduate mathematics examination},
  pdfdisplaydoctitle=true,
  bookmarksopen=true
}
\setcounter{secnumdepth}{0}
\setlength{\parindent}{0pt}
\setlength{\parskip}{0.28em}
\setlength{\emergencystretch}{3em}
\setlength{\leftmargini}{2.15em}
\setlength{\leftmarginii}{2.65em}
\setlength{\leftmarginiii}{3em}
\pagestyle{empty}
\raggedbottom
\allowdisplaybreaks
\definecolor{ProblemConcern}{HTML}{7A1F1F}
\newenvironment{problemconcern}{%
  \par\tagstructbegin{tag=Aside}\begingroup\color{ProblemConcern}
}{%
  \par\endgroup\tagstructend\nopagebreak[4]\vspace{0.12em}
}
\NewTaggingSocketPlug{block/list/label}{examproblem}{%
  \tagstructbegin{tag=Lbl}\tagstructend%
  \tagstructbegin{tag=\LBody}%
  \tagstructbegin{tag=\ProblemHeadingTag,title-o={\ProblemHeadingText}}%
  \tagmcbegin{tag=\ProblemHeadingTag,actualtext={\ProblemHeadingText}}%
  #2%
  \tagmcend%
  \tagstructend%
}
\newenvironment{examproblems}[2]{%
  \def\ProblemHeadingTag{#2}%
  \AssignTaggingSocketPlug{block/list/label}{examproblem}%
  \begin{enumerate}%
  \setcounter{enumi}{#1}%
  \renewcommand{\labelenumi}{\textbf{\theenumi.}}%
  \setlength{\topsep}{0.35em}%
  \setlength{\partopsep}{0pt}%
  \setlength{\itemsep}{0.48em}%
  \setlength{\parsep}{0.18em}%
}{\end{enumerate}}
\newenvironment{examsubparts}{%
  \AssignTaggingSocketPlug{block/list/label}{default}%
  \begin{enumerate}%
  \setlength{\topsep}{0.18em}%
  \setlength{\partopsep}{0pt}%
  \setlength{\itemsep}{0.16em}%
  \setlength{\parsep}{0.1em}%
}{\end{enumerate}}
\newenvironment{examinstructions}{%
  \par\begingroup\itshape
}{%
  \par\endgroup\vspace{0.18em}
}
\makeatletter
\renewcommand\subsection{\@startsection{subsection}{2}{0pt}%
  {-1.25ex plus -.25ex minus -.1ex}%
  {0.55ex plus .1ex}%
  {\normalfont\large\bfseries}}
\renewcommand{\@maketitle}{%
  \newpage\null\vskip -1em%
  {\footnotesize\bfseries
    \href{https://www.ufl.edu/}{University of Florida}, \href{https://math.ufl.edu/}{Department of Mathematics}\par}%
  \vskip -0.5em%
  \tagstructbegin{tag=H1,title={Analysis PhD Exam, January 1997}}%
  \tagpdfparaOff%
  \tagmcbegin{tag=H1}%
    {\LARGE\bfseries\href{https://gma.math.ufl.edu/exams/analysis-phd/}{Analysis PhD Exam}, January 1997\par}%
  \tagmcend%
  \tagpdfparaOn%
  \tagstructend%
  \vskip 0.25em%
  \tagmcbegin{artifact}%
    \hrule height 1pt%
  \tagmcend%
  \vskip 1em%
}
\makeatother
\title{Analysis PhD Exam, January 1997}
\author{}
\date{}
\begin{document}
\maketitle
\thispagestyle{empty}
\subsection{I.}
\begin{examinstructions}
State the following theorems.
\end{examinstructions}
\begin{examproblems}{0}{H3}
\def\ProblemHeadingText{Problem 1.}
\item
The monotone convergence theorem for nets in \(L^p\).
\def\ProblemHeadingText{Problem 2.}
\item
The Hahn decomposition theorem.
\def\ProblemHeadingText{Problem 3.}
\item
The Egorov theorem.
\def\ProblemHeadingText{Problem 4.}
\item
The Vitali convergence theorem.
\def\ProblemHeadingText{Problem 5.}
\item
The Radon–Nikodym theorem.
\def\ProblemHeadingText{Problem 6.}
\item
The Fubini theorem.
\def\ProblemHeadingText{Problem 7.}
\item
The representation of the dual of \(L^p\).
\def\ProblemHeadingText{Problem 8.}
\item
The Lebesgue dominated convergence theorem.
\end{examproblems}
\subsection{II}
\begin{examinstructions}
Prove one of the theorems 6 or 7.
\end{examinstructions}
\subsection{III.}
\begin{examinstructions}
Solve 5 of the following problems.
\end{examinstructions}
\begin{examinstructions}
In what follows, \((X,\Sigma,\mu)\) is a measure space.
\end{examinstructions}
\begin{examproblems}{0}{H3}
\def\ProblemHeadingText{Problem 1.}
\item
Let \((f_n)\) be a sequence of real-valued measurable functions on~\(X\), and let~\(f\) be a real measurable function on~\(X\). Prove that \((f_n)\) converges to~\(f\) in~\(\mu\)-measure, \(f_n\xrightarrow{\mu}f\), iff every subsequence \((f_{n_i})\) contains a further subsequence \((f_{n_{i_j}})\) converging to~\(f\), \(\mu\)-a.e.
\def\ProblemHeadingText{Problem 2.}
\item
Let \((f_n)\) be a sequence of positive, \(\mu\)-integrable functions such that \(\sum_n\int f_n\,d\mu<\infty\). For each \(x\in X\), denote
\[
L(x)=\sup\{n:f_n(x)>1\}.
\]
 Prove that \(L(x)<\infty\), \(\mu\)-a.e.
\def\ProblemHeadingText{Problem 3.}
\item
Let \(f:\mathbb R\to\mathbb R\) be a Lebesgue integrable function. Prove that for every \(\varepsilon>0\), there is a bounded interval~\(I\) such that \(\left|\int_E f\,dx\right|<\varepsilon\) for every measurable set~\(E\) with \(E\cap I=\varnothing\).
\def\ProblemHeadingText{Problem 4.}
\item
Let~\(F\) be a Banach space and \(\mathcal F\subseteq\Sigma\) a sub-σ-algebra. Prove the existence of the conditional expectation \(E(f\mid\mathcal F)\) for every~\(\mu\)-integrable function \(f:X\to F\).
\def\ProblemHeadingText{Problem 5.}
\item
Let \((X,\mathcal S,\mu)\) and \((Y,\mathcal T,\nu)\) be two real measure spaces. Prove the existence of the product measure \(\mu\times\nu\) and that \(|\mu\times\nu|=|\mu|\times|\nu|\).
\def\ProblemHeadingText{Problem 6.}
\item
Let~\(\mu\) be the Lebesgue measure on \(\mathbb R\), and let~\(f\) be a Lebesgue integrable function on \(\mathbb R\) such that \(\int_0^x f\,d\mu=0\) for every \(x\in\mathbb R\) (if \(x<0\), \(\int_0^x=-\int_x^0\)). What can you say about~\(f\)?
\def\ProblemHeadingText{Problem 7.}
\item
Let~\(\mu\) be the Lebesgue measure on \(\mathbb R\), and let \(f_n=-\varphi_{(n,\infty)}\), \(n=1,2,\ldots\). Show that \((f_n)\) is increasing but that the conclusion of the monotone convergence theorem is not true. Explain why.
\def\ProblemHeadingText{Problem 8.}
\item
\begin{problemconcern}
\textbf{Warning: Suspected error.} The source legibly uses \((0,1)\) for the measure space but \((0,1]\) for the binary-expansion domain; whether this endpoint difference is intentional is unclear.
\end{problemconcern}
Let~\(\mu\) be the Lebesgue measure on \((0,1)\). Each point \(x\in(0,1]\) has a binary expansion
\[
x=\sum_{i=1}^{\infty}\frac{x_i}{2^i},
\]
 where \(x_i=0\) or~\(1\). (The expansion is unique except on a countable set, which we can throw out of the space.)
\begin{examsubparts}
\item[\textnormal{(i)}]
Define \(T_i(x)=x_i\). Sketch the graphs of \(T_1,T_2,T_3\).
\item[\textnormal{(ii)}]
Let \(\mathcal F_k=\sigma(T_1,T_2,\ldots,T_k)\). Describe \(\mathcal F_k\).
\item[\textnormal{(iii)}]
Let \(f\in L^2(\mu)\). Find a function \(g_1\) which is \(\mathcal F_1\)-measurable and satisfies
\[
\int_A g_1\,d\mu=\int_A f\,d\mu\quad\text{for }A\in\mathcal F_1.
\]
 That is, find \(g_1=E(f\mid\mathcal F_1)\). Find \(g_2=E(f\mid\mathcal F_2)\).
\end{examsubparts}
\end{examproblems}
\end{document}
